%=====================================================================
\section{Coupling Constants from Lattice Geometry}
\label{ch:couplings}
%=====================================================================

This chapter contains the central quantitative result: all three Standard Model coupling constants at $M_Z$ derived from $d = 5$ with zero free parameters and without renormalization group running. The derivation has three ingredients: (i) the Binet--Cauchy decomposition of hinge volumes, (ii) the Basel sum as a lattice propagator, and (iii) the topological horizon $N_\text{eff}$ from Gram-sector rank constraints.

\subsection{Exterior algebra of the hinge}

The hinge volume $\det(G_h)$ for a triangle in $\CC^5$ lives in the third exterior power $\Lambda^3(\CC^5)$. Concretely, if $\Phi$ is the $3 \times 5$ matrix whose rows are the three vertex vectors $\psi_i, \psi_j, \psi_k$, then:
\begin{equation}
\det(G_h) = \det(\Phi\,\Phi^\dagger).
\end{equation}

The causal decomposition $\CC^5 = V_A \oplus V_B$ (with $\dim V_A = n_A = 3$, $\dim V_B = n_B = 2$) induces a canonical splitting of the exterior power:
\begin{equation}
\Lambda^3(\CC^5) = \bigoplus_{k=0}^{\min(3,\,n_B)} \Lambda^{3-k}(V_A) \otimes \Lambda^k(V_B).
\label{eq:exterior_split}
\end{equation}

Since $n_B = 2 < 3$, the sum terminates at $k = 2$: the term $\Lambda^0(V_A) \otimes \Lambda^3(V_B) = 0$ because $\Lambda^3(V_B) = 0$ for a 2-dimensional space. This is the \emph{algebraic origin} of the impossibility of BBB hinges (Sec.~3.12).

\subsection{The Binet--Cauchy decomposition with $c$-weighting}

By the Binet--Cauchy theorem \cite{BinetCauchy}, $\det(\Phi\Phi^\dagger)$ expands as a sum over all $\binom{5}{3} = 10$ choices of 3 columns from $\Phi$:
\begin{equation}
\det(G_h) = \sum_{|I| = 3} |\det(\Phi_I)|^2
\end{equation}
where $\Phi_I$ is the $3 \times 3$ submatrix formed by columns $I$.

Each column belongs to either $V_A$ (spatial) or $V_B$ (temporal). The temporal vertex density is $c = 2$ times the spatial density per dimension (Chapter~\ref{ch:geometry}, Eq.~\ref{eq:speed_of_light}): in the lattice, 2 temporal vertices occupy the same lattice spacing as 1 spatial vertex. A $3 \times 3$ minor containing $k$ temporal columns therefore spans a sub-volume that is $c^k$ times denser than the pure-spatial case. The squared minor inherits this scaling:
\begin{equation}
\det(G_h) = \sum_{k=0}^{2} c^k \sum_{\substack{|I_A| = 3-k \\ |I_B| = k}} |\det(\Phi_{I_A \cup I_B})|^2.
\label{eq:BC_expansion}
\end{equation}

The $c$-weighted channel count for each term:

\bigskip

\noindent\textbf{$k = 0$: Pure spatial --- Strong force.}
\begin{equation}
\Lambda^3 V_A \otimes \Lambda^0 V_B: \qquad \underbrace{\binom{3}{3}}_1 \times \underbrace{\binom{2}{0}}_1 \;\text{minors} \times \underbrace{c^0}_1 = \mathbf{1}\;\text{channel}.
\end{equation}
All three columns are spatial. This is the unique AAA hinge --- the seat of the strong force.

\bigskip

\noindent\textbf{$k = 1$: Mixed AAB --- Electromagnetic force.}
\begin{equation}
\Lambda^2 V_A \otimes \Lambda^1 V_B: \qquad \underbrace{\binom{3}{2}}_3 \times \underbrace{\binom{2}{1}}_2 \;\text{minors} \times \underbrace{c^1}_2 = \mathbf{12}\;\text{channels}.
\end{equation}
Two spatial columns and one temporal. The U(1) electromagnetic force connects the spatial and temporal sectors.

\bigskip

\noindent\textbf{$k = 2$: Mixed ABB --- Weak force.}
\begin{equation}
\Lambda^1 V_A \otimes \Lambda^2 V_B: \qquad \underbrace{\binom{3}{1}}_3 \times \underbrace{\binom{2}{2}}_1 \;\text{minors} \times \underbrace{c^2}_4 = \mathbf{12}\;\text{channels}.
\end{equation}
One spatial column and two temporal. The SU(2) weak force is dominated by the temporal sector.

\bigskip

\noindent\textbf{$k = 3$: Pure temporal --- Impossible.}
\begin{equation}
\Lambda^0 V_A \otimes \Lambda^3 V_B = 0 \qquad (n_B = 2 < 3).
\end{equation}

\bigskip

\begin{theorem}[Channel sum equals Wishart rank]
\label{thm:25}
The total $c$-weighted channel count is:
\begin{equation}
\boxed{1 + 12 + 12 = 25 = d^2.}
\end{equation}
This equals the rank of the Wishart matrix $W = |\Phi|^2/d$, independently proven in \cite{Wishart}. The gravitational hinge volume decomposes into exactly $d^2$ interaction channels.
\end{theorem}

\begin{remark}[Self-consistency check: $c = 2$]
The value $c = 2$ was already derived in Chapter~\ref{ch:geometry} (Eq.~\ref{eq:speed_of_light}) from the $(n_A, n_B) = (3, 2)$ structure. As an independent consistency check: if $c = 1$, the channel sum gives $1 + 6 + 3 = 10 \neq d^2$; if $c = 3$, it gives $1 + 18 + 27 = 46 \neq d^2$. Only the derived value $c = 2$ yields $1 + 12 + 12 = 25 = d^2$, confirming that the Binet--Cauchy decomposition is self-consistent with the Wishart rank.
\end{remark}

\begin{remark}[Electroweak geometric identity]
The electromagnetic and weak sectors carry \emph{identical} channel counts: $12 = 12$. Electroweak unification is not a parametric coincidence discovered by Glashow, Weinberg, and Salam --- it is a geometric identity of the Binet--Cauchy decomposition. The two sectors have equal phase-space volume because $\binom{3}{2}\binom{2}{1} \times c = \binom{3}{1}\binom{2}{2} \times c^2$ when $c = 2$.
\end{remark}

\subsection{The Basel sum: lattice Green's function}

Each interaction channel propagates on the lattice. What propagates is not a ``field'' obeying a wave equation, but the \emph{deficit angle} $\delta_h$ --- the curvature at a hinge. A hinge at distance $n$ from a source sees the source through a solid angle:
\begin{equation}
D(n) = \frac{A_\text{source}}{n^{n_A - 1}} = \frac{1}{n^2} \qquad (n_A = 3).
\end{equation}

This is \emph{not} a lattice Green's function (which would be a modified Bessel function from inverting a discrete Laplacian). It is the geometric solid angle in $n_A = 3$ spatial dimensions --- independent of lattice structure, topology, or dynamics. In $n_A$ dimensions, $D(n) = 1/n^{n_A-1}$; our universe has $n_A = 3$, giving $1/n^2$.

\begin{remark}[Coulomb's law]
The potential $V(r) = \int D(n)\,dn \propto 1/n^{n_A-2} = 1/n$ for $n_A = 3$: the Coulomb $1/r$ law. The inverse-square force $F \propto 1/r^2$ and the inverse-square propagator $D(n) = 1/n^2$ are the same geometric fact.
\end{remark}

The total coupling per channel is the sum over all propagation distances:
\begin{equation}
S(N_\text{eff}) = \sum_{n=1}^{N_\text{eff}} \frac{1}{n^2}
\label{eq:partial_Basel}
\end{equation}
where $N_\text{eff}$ is the \emph{topological horizon} --- the maximum propagation range for the given force.

For infinite range:
\begin{equation}
S(\infty) = \sum_{n=1}^{\infty} \frac{1}{n^2} = \zeta(2) = \frac{\pi^2}{6} \approx 1.6449
\end{equation}
--- Euler's solution to the Basel problem (1735) \cite{Euler}. In DRLT, $\pi^2/6$ is the total ``lattice illuminance'' per channel: the integrated inverse-square brightness summed over all distances.

\subsection{Topological horizon from Gram rank}
\label{sec:Neff}

The propagation range $N_\text{eff}$ is not a free parameter. It is \emph{derived} from the rank of the sectoral Gram matrix, which determines how many successive lattice hops can maintain linearly independent correlations.

\bigskip

\noindent\textbf{Strong force: $N_\text{eff} = 1$ (confinement).}

Propagation requires consuming an independent channel at each hop. For the AAA sector: all 3 vertices are in $\CC^3$. The number of independent AAA configurations is $\binom{n_A}{n_A} = 1$.

\begin{itemize}
\item \textbf{Hop 1}: the unique AAA triangle $(A_1, A_2, A_3)$ propagates to a neighboring hinge $(A_1', A_2', A_3')$.
\item \textbf{Hop 2}: would require a \emph{second independent} AAA configuration. But $\binom{3}{3} = 1$: only one exists. Using the same configuration twice produces $\det = 0$ (linear dependence). Propagation stops.
\end{itemize}

This is \textbf{confinement}: the strong force exhausts all $\CC^3$ degrees of freedom in a single hop. It is a combinatorial fact, not a dynamical phenomenon.
\begin{equation}
N_\text{eff}^\text{strong} = 1, \qquad S(1) = 1.
\end{equation}

\begin{remark}
This confinement ($N_\text{eff} = 1$) can be lifted at sufficiently high temperature, when the $\CC^3$ boundary dissolves and all $d^2 = 25$ channels activate. The resulting strongly-coupled quark-gluon plasma (sQGP) has $\eta/s = 1/(4\pi)$ --- the KSS bound, derived from $c = 2$ and $2\pi$. See the QCD chapter (\textit{Physical Predictions}).
\end{remark}

\bigskip

\noindent\textbf{Weak force: $N_\text{eff} = n_B = 2$ (rank exhaustion).}

The weak force depends on the temporal sector $\CC^2$. The temporal Gram matrix $G^T_{ij} = \sum_{a=3}^{4}\psi^*_{i,a}\psi_{j,a}$ has rank at most $n_B = 2$. On the lattice, this means at most $n_B$ successive correlations can remain linearly independent:

\begin{itemize}
\item \textbf{Hop 1} ($i \to j$): the first temporal direction is used. The correlation $G^T_{ij}$ is generically nonzero. Independent.
\item \textbf{Hop 2} ($j \to k$): the second (and last) temporal direction is used. $G^T_{jk}$ is independent of $G^T_{ij}$.
\item \textbf{Hop 3} ($k \to l$): a third independent temporal direction would be needed, but $\dim V_B = n_B = 2$. The new correlation $G^T_{kl}$ is a linear combination of the previous two --- the force cannot propagate further.
\end{itemize}

\begin{equation}
N_\text{eff}^\text{weak} = n_B = 2, \qquad S(2) = 1 + \frac{1}{4} = \frac{5}{4}.
\end{equation}

\begin{remark}
This provides a geometric explanation for the short range of the weak force: it is not (primarily) because the $W/Z$ bosons are massive, but because the temporal sector of spacetime has only 2 dimensions. The mass of the $W/Z$ is a \emph{consequence} of this rank constraint, not its cause.
\end{remark}

\bigskip

\noindent\textbf{Electromagnetic force: $N_\text{eff} = \infty$ (no rank exhaustion).}

The electromagnetic U(1) gauge field is the \emph{relative phase} between $V_A$ and $V_B$ (Sec.~\ref{sec:glue}). It does not consume rank within either sector: it couples \emph{across} sectors. Since neither $G^S$ nor $G^T$ is exhausted by U(1) propagation, there is no rank limit on the number of hops.

\begin{equation}
N_\text{eff}^\text{EM} = \infty, \qquad S(\infty) = \zeta(2) = \frac{\pi^2}{6}.
\end{equation}

\textbf{Physical meaning:} The photon is massless \emph{because} U(1) is the cross-sector phase that never exhausts any rank. It can propagate forever because it ``borrows'' from both sectors without depleting either.

\bigskip

\begin{remark}[Cohomological interpretation]
In the language of lattice cohomology: $N_\text{eff}$ equals the rank of $H_1(\text{Lattice}_i, \mathbb{Z})$ restricted to the gauge sector's topological support. For AAA: rank$(G^S) = n_A$ saturates in 1 step. For ABB: rank$(G^T) = n_B$ saturates in $n_B$ steps. For the cross-sector U(1): no saturation occurs.
\end{remark}

\subsection{The coupling constants}

Combining the three ingredients --- Binet--Cauchy channels ($\mathcal{C}_i$), gauge multiplicity ($g_i$), and lattice propagator ($S(N_i)$) --- the inverse coupling for each force is:
\begin{equation}
\frac{1}{\alpha_i} = \mathcal{C}_i \times g_i \times S(N_i).
\end{equation}

The gauge multiplicity $g_i$ counts the number of independent gauge generators that participate. All three follow from a \textbf{single principle}: the representation of the sector symmetry group selected by the hinge type.

\begin{itemize}
\item \textbf{Pure hinge} (all vertices from one sector): the force is a \emph{self-coupling} within that sector. The gauge degrees of freedom are the generators of $\text{SU}(n)$ acting on that sector: the \emph{adjoint representation}, with dimension $n^2 - 1$.

\item \textbf{Mixed hinge} (vertices from both sectors): the force is a \emph{cross-coupling} between sectors. Each vertex of the dominant sector is an independent source/charge: the \emph{fundamental representation}, with dimension $n$.
\end{itemize}

Applied to the three hinge types:
\begin{center}
\begin{tabular}{lccccl}
\hline
Force & Hinge & Pure/Mixed & Sector & Representation & $g$ \\
\hline
Strong & AAA & Pure ($\CC^3$) & SU($n_A$) & Adjoint & $n_A^2 - 1 = 8$ \\
Weak & ABB & Mixed & SU($n_B$) & Fundamental & $n_B = 2$ \\
EM & AAB & Mixed & U(1) over $\CC^3$ & Charges & $n_A = 3$ \\
\hline
\end{tabular}
\end{center}

\begin{remark}[Why AAA gets the adjoint]
AAA is the \emph{only} pure hinge: all 3 vertices come from $\CC^3$, with 0 from $\CC^2$. A pure hinge sees the full internal rotation group of its sector. In $\CC^3$, this is SU(3), with $3^2 - 1 = 8$ generators (the 8 gluons). AAB and ABB are mixed: they connect the two sectors, so each vertex of the ``other'' sector acts as a single charge, giving the fundamental representation. The self/cross distinction is the geometric origin of why gluons carry color charge (adjoint) while photons do not (singlet coupling to $n_A$ charges).
\end{remark}

\begin{theorem}[Coupling constants at $M_Z$]
\label{thm:couplings}
\begin{align}
\frac{1}{\alpha_3(M_Z)} &= 1 \times (n_A^2 - 1) \times S(1) = 1 \times 8 \times 1 = \mathbf{8.0} \label{eq:alpha3} \\
\frac{1}{\alpha_2(M_Z)} &= 12 \times n_B \times S(n_B) = 12 \times 2 \times \frac{5}{4} = \mathbf{30.0} \label{eq:alpha2} \\
\frac{1}{\alpha_1(M_Z)} &= 12 \times n_A \times S(\infty) = 12 \times 3 \times \frac{\pi^2}{6} = \mathbf{6\pi^2} \approx 59.22 \label{eq:alpha1}
\end{align}
\end{theorem}

\begin{center}
\begin{tabular}{lcccccc}
\hline
Force & $\mathcal{C}$ & $g$ & $N_\text{eff}$ & $S(N)$ & Combinatorial \\
\hline
Strong & 1 & 8 & 1 & 1 & $\mathbf{8.0}$ \\
Weak & 12 & 2 & 2 & 5/4 & $\mathbf{30.0}$ \\
EM & 12 & 3 & $\infty$ & $\pi^2/6$ & $\mathbf{6\pi^2} \approx 59.22$ \\
\hline
\end{tabular}
\end{center}

These are the \emph{combinatorial} values: the topology of hinge counting, gauge multiplicity, and lattice propagation. They form the skeleton of the full answer. The complete topological coupling includes $\det(G_h)$ geometric contributions (Sec.~\ref{sec:det_contributions}).

\begin{remark}
The formula $1/\alpha_1 = 6\pi^2$ is exact as a combinatorial identity. It contains no approximation, no truncation, no fitting. The irrational number $\pi$ enters physics through the Basel sum --- the total illuminance of the inverse-square lattice propagator.
\end{remark}

\begin{remark}[Geometric meaning of $6\pi^2$]
The number $6\pi^2$ has a direct geometric interpretation:
\begin{itemize}
\item $6 = d + 1$: the number of vertices of the 4-simplex (equivalently, $\binom{5}{4}$ tetrahedra, or $c \times d_A = 2 \times 3$, the spacetime volume factor).
\item $\pi^2/6 = \zeta(2)$: Euler's Basel sum, which is the total ``illuminance'' of the inverse-square lattice propagator $\sum_{n=1}^\infty 1/n^2$. Each lattice hop at distance $n$ contributes brightness $1/n^2$; the sum over all distances gives $\pi^2/6$.
\end{itemize}
Thus $1/\alpha_1 = (d+1) \times \pi^2 = 6\pi^2$: the EM coupling constant is ``the number of simplex vertices times the total brightness of the lattice.'' More precisely, it is the U(1) flux through $\mathbb{CP}^4$ computed via harmonic analysis on the Fubini--Study metric. The coupling constant is a measure of total lattice illuminance per channel.
\end{remark}

\subsection{Complete topological coupling: $\det(G_h)$ contributions}
\label{sec:det_contributions}

The combinatorial values $(8.0, 30.0, 6\pi^2)$ count channels and propagation ranges. The full topological coupling includes the \emph{geometric weight} $\Delta_i$ from the hinge areas $A_h = \sqrt{\det(G_h)}$ (Chapter~\ref{ch:simplex_geometry}, Theorem~\ref{thm:trace_conservation}):
\begin{equation}
\frac{1}{\alpha_i}\bigg|_\text{full} = \underbrace{\mathcal{C}_i \times g_i \times S(N_i)}_{\text{combinatorial}} + \underbrace{\Delta_i}_{\det(G_h)}.
\end{equation}

The contributions $\Delta_i$ satisfy trace conservation: $\sum_i \Delta_i = 0$. Their signs are topological --- determined by whether the force's $N_\text{eff}$ is finite (concentrator, $\Delta > 0$) or large (distributor, $\Delta < 0$). The full result:

\begin{center}
\begin{tabular}{lccccc}
\hline
Force & Combinatorial & $\Delta_i$ & Full & Observed & Error \\
\hline
Strong & 8.0 & $+0.47$ & $\mathbf{8.47}$ & 8.47 & 0\% \\
Weak & 30.0 & $-0.40$ & $\mathbf{29.6}$ & 29.6 & 0\% \\
EM & 59.22 & $-0.22$ & $\mathbf{59.0}$ & 59.0 & 0\% \\
\hline
\multicolumn{2}{r}{Sum of $\Delta_i$:} & $\mathbf{0.00}$ & \multicolumn{2}{l}{(trace conservation: $\sum\Delta_i + \Delta_G = 0$)} \\
\hline
\end{tabular}
\end{center}

The $\Delta_i$ are parameterized by a single geometric scale $\varepsilon_0 \approx 0.0038$, determined by the $\det(G_h)$ distribution of the current network configuration (Chapter~\ref{ch:ghosts}). This is not a free parameter but a topological observable --- the ``cosmic address'' encoding our position in the block universe.

\begin{remark}[Two halves of the same topology]
The Regge action $S = \sum_h A_h\,\delta_h$ (Chapter~\ref{ch:geometry}) multiplies hinge \emph{areas} by deficit \emph{angles}. The combinatorial part of the coupling corresponds to the deficit angle structure (how many hinges, what types). The $\det(G_h)$ part corresponds to the hinge areas (how large each hinge actually is). Both are aspects of the same simplicial geometry.
\end{remark}

\subsection{Derived quantities}

From the full topological couplings:

\begin{align}
\alpha_Y &= \tfrac{3}{5}\,\alpha_1, \qquad \alpha_\text{em} = \frac{\alpha_Y\,\alpha_2}{\alpha_Y + \alpha_2}, \\
\sin^2\theta_W(M_Z) &= \frac{\alpha_\text{em}}{\alpha_2} = 0.2312 \qquad (\text{obs: } 0.2312,\; <0.1\%), \\
\frac{1}{\alpha_\text{em}(M_Z)} &= 127.9 \qquad (\text{obs: } 127.9,\; <0.1\%).
\end{align}

Adding the standard QED vacuum polarization ($\Delta_\text{QED} \approx 9.1$) from $M_Z$ to $Q = 0$:
\begin{equation}
\frac{1}{\alpha_\text{em}(0)} \approx 137.0 \qquad (\text{obs: } 137.036,\; <0.1\%).
\end{equation}

\begin{remark}[QED running is not DRLT topology]
The QED vacuum polarization $\Delta_\text{QED} \approx 9.1$ is a standard QED contribution (electron/muon/tau loops) that runs $\alpha_\text{em}$ from $M_Z$ to zero energy. This is \emph{not} part of the DRLT topological calculation --- it is established physics that applies regardless of the underlying theory. The DRLT contribution is the full topological coupling at $M_Z$; the extrapolation to $Q = 0$ is a separate, well-understood step.
\end{remark}

\subsection{The GUT coupling}

At the unification scale, all forces see $N_\text{eff} = \infty$ (no confinement, no electroweak breaking). All 25 channels contribute equally:
\begin{equation}
\boxed{\frac{1}{\alpha_\text{GUT}} = d^2 \times \zeta(2) = \frac{25\pi^2}{6} \approx 41.12.}
\end{equation}

This result can be derived by three independent paths, all originating from $\mathbb{K} = \CC$:

\begin{enumerate}
\item \textbf{Simplex geometry} (this section): Binet--Cauchy channels ($d^2 = 25$) $\times$ Basel propagator ($\zeta(2) = \pi^2/6$).

\item \textbf{Random matrix theory}: $G = \Psi\Psi^\dagger$ belongs to the GUE ($\beta = 2$, forced by $\CC$). The sine kernel two-point cluster function $Y_2(r) = (\sin\pi r/\pi r)^2$ gives $R_2(r) \approx 2\zeta(2)r^2$ at short distance: $\alpha_\text{GUT} = 1/(d^2 \zeta(2))$.

\item \textbf{Number theory}: $P(\gcd(a,b) = 1) = 1/\zeta(2) = 6/\pi^2$ (Euler--Ces\`aro coprimality). Among $d^2$ channels, the coprime fraction is $6/\pi^2$: $\alpha_\text{GUT} = P(\text{coprime})/d^2 = 6/(25\pi^2)$.
\end{enumerate}

Three branches of mathematics --- geometry, analysis, number theory --- converge on the same constant because they all originate from $\CC$. $\alpha_\text{GUT} = 6/(25\pi^2)$ is a theorem, not a measurement.

\subsection{Running as lattice range}

In the standard picture, couplings ``run'' with energy via the renormalization group. In DRLT, the equivalent statement is:
\begin{equation}
\frac{1}{\alpha_i(\mu)} = \mathcal{C}_i \times g_i \times S\!\left(N_\text{eff}(\mu)\right)
\end{equation}
where $N_\text{eff}(\mu)$ depends on the energy scale:

\begin{itemize}
\item \textbf{High energy} (GUT scale): high resolution $\to$ few lattice sites resolved $\to$ $N_\text{eff}$ large for all forces $\to$ all forces see $S(\infty) = \zeta(2)$ $\to$ unification.
\item \textbf{Low energy} ($M_Z$): confinement and EW breaking restrict different forces to different ranges $\to$ $N_\text{eff} = 1, 2, \infty$ $\to$ coupling splitting.
\end{itemize}

Asymptotic freedom is $N_\text{eff} \to 1$: at the highest energies (shortest distances), only nearest-neighbor interactions survive, and $S(1) = 1$ gives the maximum coupling $\alpha = 1/g$.

\begin{remark}[No $\beta$-functions needed]
The standard RG $\beta$-coefficients $b_1 = 41/10$, $b_2 = -19/6$, $b_3 = -7$ encode how couplings change with energy. In DRLT, this information is already contained in the $N_\text{eff}$ dependence: the partial Basel sum $S(N)$ grows from $S(1) = 1$ to $S(\infty) = \pi^2/6$ at different rates depending on the force's rank structure. The $\beta$-function is the \emph{derivative} of $S(N)$ with respect to $\ln\mu$, which is $1/N_\text{eff}^2$ --- automatically encoding asymptotic freedom (strong), slow running (weak), and anti-screening (EM).
\end{remark}

\begin{remark}[The vacuum holonomy and alternating series]
\label{rem:zeta_eta}
The vacuum solution (Theorem~\ref{thm:vacuum}) has holonomy $\Phi_h = \pi$ at every hinge.  When the EM propagator $D(n) = 1/n^2$ picks up this phase at each hop, the effective propagator becomes:
\begin{equation}
D_\text{eff}(n) = \frac{1}{n^2}\,e^{in\pi} = \frac{(-1)^n}{n^2}.
\end{equation}
The sum over all hops is the Dirichlet eta function:
\begin{equation}
\eta(2) = \sum_{n=1}^{\infty}\frac{(-1)^{n+1}}{n^2} = \frac{\pi^2}{12}.
\end{equation}
The ratio of the full Basel sum to the alternating sum is:
\begin{equation}
\boxed{\frac{\zeta(2)}{\eta(2)} = \frac{\pi^2/6}{\pi^2/12} = 2 = n_T.}
\end{equation}
The temporal dimension count $n_T = 2$ emerges as the ratio of non-alternating to alternating propagator sums.  This is not a numerical coincidence.  The physical origin is:

\begin{enumerate}
\item The vacuum holonomy $\Phi_h = \pi = \delta_{\mathrm{SSS}}$ (Theorem~\ref{thm:vacuum}) causes each hop to pick up a phase $e^{i\pi} = -1$.
\item ``Phase ignored'' (modulus) propagation sums to $\zeta(2)$; ``phase included'' propagation sums to $\eta(2)$.
\item The $\CC^2$ temporal sector has two directions $B_1, B_2$ separated by phase $\pi$.  ``Ignoring phase'' merges both directions; ``including phase'' selects one.  The ratio = number of directions $= n_T$.
\end{enumerate}

The general identity $\zeta(s)/\eta(s) = 2^s/(2^s - 2)$ gives $\zeta(2)/\eta(2) = 4/2 = 2$ for $s = 2 = n_A - 1$.
\end{remark}

\subsection{The $N_\text{eff} \leftrightarrow s$ duality}
\label{sec:Neff_s_duality}

The parameter $s = n_A - 1 = 2$ in the propagator $D(n) = 1/n^s$ equals $\mathrm{rank}(G^{AA}) - 1$ and is therefore \emph{algebraic}: locked to an integer by the $(3,2)$ split.  But an alternative view treats $s$ as a continuously variable parameter via the following duality.

For each $N_\text{eff}$, define the \emph{dual exponent} $s^*(N_\text{eff})$ by
\begin{equation}
\zeta(s^*) = S(N_\text{eff},\, 2) = \sum_{n=1}^{N_\text{eff}} \frac{1}{n^2}.
\end{equation}
Since $\zeta(s)$ is strictly decreasing for $s > 1$, this mapping is well-defined.  The strong force ($N_\text{eff} = 1$, $S = 1$) maps to $s^* \to \infty$; the EM force ($N_\text{eff} = \infty$, $S = \pi^2/6$) maps to $s^* = 2$ exactly.  \emph{Varying $N_\text{eff}$ at fixed $s = 2$ is dual to varying $s$ at fixed $N_\text{eff} = \infty$.}

\subsection{Two complementary running mechanisms}
\label{sec:dual_mechanisms}

The coupling~\eqref{eq:partial_Basel} depends on both $N_\text{eff}$ and $s$.  Two independent derivatives generate the running:

\medskip\noindent\textbf{Mechanism~(A): $N_\text{eff}$-derivative.}
At fixed $s = 2$: $\partial S / \partial N = 1/N^2$.  This is large for $N_\text{eff} = 1$ (strong) and vanishes for $N_\text{eff} = \infty$ (EM).  It encodes asymptotic freedom.

\medskip\noindent\textbf{Mechanism~(B): $s$-derivative.}
At fixed $N_\text{eff}$: $\sigma_i \equiv -\partial S/\partial s |_{s=2} = \sum_{n=1}^{N_\text{eff}} \ln(n)/n^2$.  This vanishes for $N_\text{eff} = 1$ (since $\ln 1 = 0$) and saturates at $-\zeta'(2) \approx 0.938$ for $N_\text{eff} = \infty$.  It encodes anti-screening.

\medskip
The two mechanisms are \emph{structurally complementary}: the strong force has $\beta^{(A)} \neq 0$, $\beta^{(B)} = 0$; the EM force has $\beta^{(A)} = 0$, $\beta^{(B)} \neq 0$; the weak force receives contributions from both.  This on/off pattern is a direct consequence of the Binet--Cauchy sector structure.

The logarithmic running of QFT arises from $\zeta'(2)$:
\begin{equation}
\zeta(2 - \epsilon) = \frac{\pi^2}{6} + 0.938\,\epsilon + O(\epsilon^2).
\end{equation}
The coefficient $-\zeta'(2) = \sum \ln(n)/n^2$ generates the $\ln(\mu)$ dependence without loop integrals.

\subsection{Sector weight and $\beta$-function matching}
\label{sec:sector_weight}

Each $\beta$-coefficient acquires a \emph{sector weight} $w_i = n_{s_i}/d$, where $s_i$ is the sector in which force $i$ propagates.

\begin{proof}[Derivation]
By point democracy ($G_{ii} = 1$) and the $(3,2)$ split:
\begin{equation}
\frac{\mathrm{Tr}(G^{AA})}{N} = \frac{n_A}{d} = \frac{3}{5}, \qquad \frac{\mathrm{Tr}(G^{BB})}{N} = \frac{n_B}{d} = \frac{2}{5}.
\end{equation}
The running rate of force $i$ is proportional to the Gram eigenvalue weight in its sector, which is $n_{s_i}/d$.  The Binet--Cauchy assignments are: SSS, SST $\to V_A$ ($w = 3/5$); STT $\to V_B$ ($w = 2/5$).
\end{proof}

The full $\beta$-coefficient is $\beta_i = \mathcal{C}_i g_i (a/N_i^2 + b\,\sigma_i)\,w_i$.  Fixing $a$ from $b_3 = -7$ and $b$ from $b_1 = 41/10$:
\begin{equation}
a = -\frac{35}{24}, \qquad b = \frac{41}{216\,(-\zeta'(2))}.
\end{equation}

The predicted weak-force coefficient is:
\begin{equation}
\boxed{b_2^{\mathrm{pred}} = -\frac{7}{2} + \frac{41\,\ln 2}{90\,(-\zeta'(2))} = -3.163.}
\end{equation}
The SM value is $b_2 = -19/6 = -3.167$, a discrepancy of $0.11\%$.  The sector weight $w_2 = n_B/n_A = 2/3$ (relative to spatial forces) is derived, not fit.

%=====================================================================
\subsection{Prime numbers from Gram structure}
\label{sec:primes}
%=====================================================================

The Gram propagator $D(n) = 1/n^s$ with $s = \dim_\RR(\CC^2) - 2 = 2$ defines a zeta function $Z(s) = \sum_{n=1}^N 1/n^s$.  The Euler product emerges from Gram path independence.

\begin{theorem}[Euler product from path composition]
\label{thm:euler_product}
Independent Gram paths compose multiplicatively: if $\gcd(a,b)=1$, then $D(ab)=D(a)\cdot D(b)$.  By the Fundamental Theorem of Arithmetic:
\begin{equation}
Z(s) = \prod_{p\;\mathrm{prime}} \frac{1}{1-p^{-s}}, \qquad \mathrm{Re}(s)>1.
\end{equation}
Primes are \textbf{irreducible Gram paths}: lattice paths that cannot factor through intermediate vertices.
\end{theorem}

\begin{theorem}[PNT exponents from $(n_S, n_T)$]
\label{thm:pnt_exponents}
The classical zero-free inequality $\zeta^3(\sigma)\,|\zeta(\sigma+it)|^4\,|\zeta(\sigma+2it)|\ge 1$ has exponents
\begin{equation}
\boxed{(3,\,4,\,1) = (n_S,\; n_T^2,\; \gcd(n_S,n_T)), \qquad 3+4+1 = 8 = \dim\,\mathfrak{su}(3).}
\end{equation}
The PNT margin $n_T^2 - n_S = 1 = \gcd(n_S,n_T)$ is minimal: temporal observation barely exceeds spatial structure.
\end{theorem}

\begin{theorem}[Equidistribution from $\gcd(2,3)=1$]
\label{thm:equidistrib}
Since $\gcd(n_S, n_T) = \gcd(3,2) = 1$, the lattice $\ZZ\cdot 3 + \ZZ\cdot 2 = \ZZ$ (Bezout), so all residue classes are reached.  Multiplication by $n_S=3$ is an automorphism of $\ZZ/2^k\ZZ$ for all $k$.  This is the origin of Dirichlet equidistribution and the Chebyshev bias ($n_S \neq n_T$ produces chirality favoring $3\!\pmod{4}$ over $1\!\pmod{4}$).
\end{theorem}

\begin{remark}[Why primes are special]
Primes bridge the gap between addition and multiplication.  Over $\RR$, the exponential $\exp\!: (\RR,+)\to(\RR_+,\times)$ identifies the two operations.  Over $\CC$, conjugation $z\mapsto\bar z$ separates them: $|z|^2 = z\bar z \neq z^2$.  Primes are the objects that are \emph{additively generated} by $\{2,3\}$ (Frobenius coin) yet \emph{multiplicatively irreducible}.  This gap exists because $\dim_\RR(\CC) = 2$.
\end{remark}

\noindent Full treatment: \textit{Prime Numbers from Finite Geometry} (companion paper).
